\documentclass[11pt]{article}
\usepackage{fontspec}
\setmainfont{Times New Roman}
\setsansfont{Arial}
\setmonofont{Consolas}
\usepackage{amsmath,amssymb,mathtools}
\usepackage{geometry}
\usepackage{microtype}
\geometry{a4paper,margin=2.28cm}
\setlength{\parindent}{1.15em}
\setlength{\parskip}{0pt}
\makeatletter
\renewcommand\footnotesize{\@setfontsize\footnotesize{7.5}{8.5}\selectfont}
\makeatother
\providecommand{\ma}{\mathfrak{a}}
\providecommand{\mb}{\mathfrak{b}}
\providecommand{\mc}{\mathfrak{c}}
\providecommand{\mA}{\mathfrak{A}}
\providecommand{\mZ}{\mathfrak{Z}}
\providecommand{\mo}{\mathfrak{o}}
\providecommand{\mr}{\mathfrak{r}}
\emergencystretch=120em
\allowdisplaybreaks
\sloppy
\hbadness=10000
\vbadness=10000
\hfuzz=2pt
\vfuzz=10pt
\raggedbottom
\begin{document}
% Unit: Paper34 section11 v001. Direct Interslavic Latin translation from audited German source lines 661--730, segments P34-S0129--P34-S0136. Source scan pages 20--22 retained as witnesses; Section12 heading P34-S0137 is intentionally reserved for Section12. Zenodo record 20822156 was rechecked for this unit at 2026-06-24T22:09Z UTC.
\section*{34. Hiperkompleksne veličiny i teorija predstavjenj}
\emph{Prodolženje: glava II, §§8--14.}

\section*{II. poglavje. Nekomutativna teorija idealov}

\subsection*{§ 11. Centar}

Centar $\mZ$ kolca $\mo$ jest množstvo vsih $z$, ktore komutujut so vsimi elementami kolca ($za=az$ za vse $a$). $\mZ$ jest komutativno kolco.

Vsaky jednostranny ideal $\ma$ v $\mo$ imaje v $\mZ$ ``ideal suženja'' $\ma\cap\mZ$. Bo, poneže i $\ma$ i $\mZ$ sut $\mZ$-moduly, $\ma\cap\mZ$ takože jest $\mZ$-modul.

Vsaky ideal $\mA$ v $\mZ$ imaje ideal razširenja: ideal $\ma=(\mA,\mo\mA)$, porodženy od $\mA$ v $\mo$. On jest dvustranny, poneže
\[
        \ma\mo=(\mo\mA,\mA)\mo=(\mo\mA\mo,\mA\mo)
        =(\mo^2\mA,\mo\mA)\subseteq\mo\mA\subseteq\ma.
\]
Ako se v $\mo$ predpostavja jediničny element, togda možno pisati $\ma=\mo\mA$. V tom slučaju važi slědujuči teorem:

\emph{Vsakomu razloženju $\mo$ na dvustranne idealy vzaimno-jednoznačno odgovarja razloženje centra: iz}
\begin{equation}
        \mo=\ma_1+\cdots+\ma_n,
        \qquad \mA_i=\ma_i\cap\mZ                         \tag{1}
\end{equation}
\emph{slěduje}
\begin{equation}
        \mZ=\mA_1+\cdots+\mA_n,
        \qquad \mo\mA_i=\ma_i,                              \tag{2}
\end{equation}
\emph{i obratno.}

\emph{Dokaz.} Nehaj $z$ bude centralny element, i nehaj
\begin{equation}
        z=z_1+\cdots+z_n                                    \tag{3}
\end{equation}
bude jego razloženje po (1). Togda $z_i$ sut centralne elementy; bo
\[
        za=z_1a+\cdots+z_na=az=az_1+\cdots+az_n;
\]
zaradi direktnosti sumy i poneže $\ma_i$ jest dvustranny,
\[
        z_i a=az_i.
\]
Zato $z_i$ leži v $\ma_i\cap\mZ=\mA_i$; zato (3) davaje tvrđeno razloženje sumy. Osobito, $e=\sum e_i$; zato $e_i$ sut centralne elementy. Nakonec
\[
        z=ze_1+\cdots+ze_n,
\]
zato
\[
        \mA_i=\mZ e_i,
        \qquad
        \mo\mA_i=\mo\mZ e_i=\mo e_i=\ma_i.
\]

Obratno, ako $\mZ=\mA_1+\cdots+\mA_n$ jest razloženje centra, togda po §9
\[
        \mA_i=\mZ e_i,
        \qquad e_i^2=e_i,
        \qquad e_i e_k=0\quad(i\ne k).
\]
Zato po relacijah ortogonalnosti \hbox{\(\S 9\)}
\[
\begin{aligned}
        \mo&=\mo e=\mo e_1+\cdots+\mo e_n
        =\mo\mZ e_1+\cdots+\mo\mZ e_n\\
        &=\mo\mA_1+\cdots+\mo\mA_n
        =\ma_1+\cdots+\ma_n .
\end{aligned}
\]
Ako teper postavimo $\ma_i\cap\mZ=\mA_i'$, togda iz prvoj česti tvrđenja znajemo, že
\[
        \mZ=\mA_1'+\cdots+\mA_n'\quad\hbox{direktno}.
\]
Ale $\mA_i'\supseteq\mA_i$, zato $\mA_i'=\mA_i$ po §4, 2.

\emph{Poslědok.} $\ma_i$ i $\mA_i$ sut jednočasno dvustranno direktno razložime ili ne.
\end{document}
